\documentclass[12pt,a4paper]{article}
\usepackage[T1]{fontenc}
\usepackage[utf8]{inputenc}
\usepackage[brazil]{babel}
\usepackage{lmodern}
\usepackage{microtype}
\usepackage{amsmath,amssymb}
\usepackage{geometry}
\geometry{left=3cm,right=2cm,top=2.6cm,bottom=2.4cm}
\usepackage{booktabs,tabularx,array}
\newcolumntype{Y}{>{\raggedright\arraybackslash}X}
\newcolumntype{P}[1]{>{\raggedright\arraybackslash}p{#1}}
\usepackage{caption}
\usepackage{placeins}
\usepackage{xurl}
\usepackage[hidelinks]{hyperref}
\usepackage{setspace}
\setstretch{1.1}
\setlength{\parindent}{1.25cm}
\setlength{\parskip}{0.25em}
\setlength{\tabcolsep}{6pt}
\renewcommand{\arraystretch}{1.1}
\captionsetup{font=small,labelfont=bf,justification=raggedright,singlelinecheck=false,skip=5pt}
\emergencystretch=2em
\clubpenalty=10000
\widowpenalty=10000
\displaywidowpenalty=10000
\raggedbottom
\hypersetup{pdftitle={Detecção de erros em telemetria sob ruído real},pdfauthor={Enzo Aurélio}}
\begin{document}
\hypersetup{pageanchor=false}
\begin{titlepage}
\centering
{\large\bfseries IFMT -- OCTAYDE\par}
\vspace{1.8cm}
{\large Enzo Aurélio\par}
\vfill
{\LARGE\bfseries Detecção de erros em telemetria\\[0.3em]sob ruído real\par}
\vspace{0.6cm}
{\large Estudo experimental de SUM8 e CRC8 com ESP32\par}
\vspace{2cm}
\begin{flushright}
\begin{minipage}{0.58\textwidth}
Relatório apresentado à disciplina TCD -- Transmissão e Comunicação de Dados.\par
\vspace{0.8em}
Professor: Ruy Oliveira.
\end{minipage}
\end{flushright}
\vfill
{\small Repositório do trabalho:\par\url{https://github.com/enzogomes2006/esp32-deteccao-erros-telemetria}\par}
\vspace{0.8cm}
Outubro de 2026
\end{titlepage}
\hypersetup{pageanchor=true}
\pagenumbering{roman}
\section*{Resumo}
Este trabalho compara a detecção de erros por soma modular de oito bits (SUM8) e CRC-8/SMBUS em um enlace UART entre dois ESP32, perturbado por um terceiro microcontrolador. O injetor aciona um transistor NPN para forçar bits do payload a zero e registra o efeito efetivamente observado. Foram realizadas sete rodadas válidas de 1000 quadros: quatro com SUM8, incluindo o controle sem ruído, e três com CRC8. O receptor executa a checagem antes de consultar uma cópia de depuração, utilizada exclusivamente para auditar a integridade do payload. SUM8 apresentou quatro erros não detectados no modo de dois alvos, correspondentes a 0,5083\% dos 787 quadros com erro dessa rodada. CRC8 apresentou zero erros não detectados nos modos ensaiados. O tempo de verificação registrado na série foi de 1 \ensuremath{\mu}s para SUM8 e 3 \ensuremath{\mu}s para CRC8. Em um ensaio complementar de interrupção da recepção, o RX completo retomou a aceitação 153,028 ms após a primeira detecção de perda. Os resultados evidenciam uma limitação da soma modular e a maior cobertura do CRC para as classes de erro utilizadas, sem generalizar a detecção para qualquer padrão de corrupção.

\textbf{Palavras-chave:} ESP32; comunicação UART; checksum; CRC8; detecção de erros.

\clearpage\tableofcontents\clearpage\pagenumbering{arabic}
\section{Introdução}
Sistemas de telemetria dependem da identificação de alterações nos dados transmitidos. Um código de detecção acrescenta redundância ao quadro e permite ao receptor decidir sobre sua aceitação. Entretanto, a aprovação na checagem não garante ausência de erro: determinadas alterações preservam o resíduo esperado.

O objetivo deste estudo foi implementar um protocolo serial, produzir perturbações reais na linha e comparar SUM8 e CRC8 quanto à ocorrência de erros não detectados. Também foram avaliados o custo de cálculo e verificação, o overhead e a recuperação após ausência temporária de recepção. A análise distingue a decisão autônoma do enlace da auditoria posterior dos dados.

\FloatBarrier
\section{Fundamentação teórica}
\subsection{Soma modular de oito bits}
O SUM8 reduz a soma dos bytes protegidos ao resto da divisão por 256. Neste protocolo, o transmissor anexa o complemento modular da soma. Para os sete bytes anteriores ao detector, o valor transmitido é:

\begin{equation}\mathrm{CHK}=\left(-\sum_{i=0}^{6}b_i\right)\bmod256.\end{equation}
O receptor soma o corpo completo, incluindo CHK, e aceita apenas se o resto for zero. Essa operação detecta uma alteração isolada de um bit, pois sua contribuição em valor absoluto é uma potência de dois entre 1 e 128. Por outro lado, alterações cuja soma seja múltipla de 256 não modificam o resíduo.

Como o circuito força apenas o nível baixo, apagar o bit mais significativo de dois bytes distintos reduz a soma em 128 + 128 = 256. Esse padrão mantém o resíduo zero e constitui uma colisão do SUM8. A aceitação é indevida quando o payload recebido difere do original.

\subsection{CRC-8/SMBUS}
O CRC opera sobre bits com aritmética polinomial em \(\mathrm{GF}(2)\), usando operações XOR. Foi adotado o polinômio \(x^{8}+x^{2}+x+1\), identificado no SMBus por 1 0000 0111 (SMBUS, s.d.). A implementação utiliza 0x07 sem o termo de grau oito, estado inicial 0x00, sem reflexão e com xorout=0x00.

Os sete bytes anteriores ao CRC representam 56 bits de dados. O catálogo de Koopman e suas notas indicam distância de Hamming 4 para o polinômio completo 0x107 em comprimentos de dados até 119 bits. Assim, para o bloco empregado, erros de um e de dois bits são detectáveis no modelo de substituição de bits (KOOPMAN, 2024a; 2024b).

Qualquer alteração não nula limitada a um único byte também tem extensão de até oito bits e não é divisível por esse gerador de grau oito. Essa propriedade é relevante para o modo de rajada utilizado. A cobertura não equivale à detecção de toda corrupção possível, nem abrange automaticamente perda de bytes ou falhas de enquadramento.

\FloatBarrier
\section{Especificação do protocolo}
O enlace de dados é unidirecional, do TX para o RX, em UART2 a 115200 baud, oito bits de dados, sem paridade e um bit de parada (8N1). Não há ACK nem retransmissão. A estrutura permanece igual nas rodadas SUM8 e CRC8; somente a técnica de cálculo de CHK é alterada.

\begin{quote}\small\ttfamily\raggedright A5 5A C3 3C | 07 seqL seqH p0 p1 p2 p3 CHK\end{quote}

\par\smallskip\noindent\begin{minipage}{\textwidth}\centering\small
\captionof{table}{Campos do quadro de dados}
\begin{tabularx}{\textwidth}{@{}l l Y@{}}\toprule
\textbf{Campo} & \textbf{Tamanho} & \textbf{Descrição}\\\midrule
Marcador & 4 bytes & A5 5A C3 3C; identifica o início e não entra na checagem\\
LEN & 1 byte & Valor 07; conta os sete bytes posteriores a LEN\\
Sequência & 2 bytes & Inteiro sem sinal de 16 bits, byte menos significativo primeiro\\
Payload & 4 bytes & Dados sintéticos de telemetria\\
CHK & 1 byte & SUM8 ou CRC8; uma técnica por rodada\\
\bottomrule\end{tabularx}
\par\smallskip{\raggedright\footnotesize Fonte: Elaboração própria, a partir da implementação do protocolo (2026).\par}\end{minipage}\par\smallskip

A checagem protege LEN, sequência e payload. No TX, o SUM8 é calculado como complemento modular da soma dos sete bytes; o CRC8 é calculado sobre esses mesmos bytes, na ordem do corpo e com processamento MSB primeiro dentro de cada byte. No RX, os oito bytes são processados e o resíduo esperado é zero. A ordem de bits do cálculo CRC não altera a transmissão LSB primeiro de cada byte pela UART.

A delimitação por comprimento fixo permite qualquer valor no payload, sem escape de bytes. Ao reconhecer o marcador, o RX coleta oito bytes e exige LEN=7. Um marcador dentro desse corpo é tratado como dado. Concluído ou descartado o corpo, a busca de marcadores recomeça no próximo byte processado.

\subsection{Referência de depuração}
Após o quadro, o TX envia D5 6A B4 4B seguido dos oito bytes originais do corpo. Essa referência contém a sequência para associação e permite validar a própria referência antes de comparar o payload. Sua leitura ocorre depois da decisão ACEITO ou DESCARTADO e não modifica o controle de sequência.

O TX insere pausas de 10 ms após o marcador de dados, 3 ms após cada byte do corpo e 10 ms antes da cópia. Essas pausas mantêm a linha em estado ocioso alto e permitem o processamento do injetor; a velocidade de transmissão dos bytes continua em 115200 baud.

\FloatBarrier
\subsection{Controle de recepção e sequência}
A rodada é delimitada por uma janela de 10 ou 1000 posições, com sequência inicial zero nos comandos de operação. O RX calcula o resíduo antes de consultar a referência. Para uma sequência admissível, registra posições anteriores ainda não observadas como perdidas, marca a posição atual como aceita ou descartada e passa a esperar o número seguinte. Um quadro descartado após checagem não é contado novamente como perdido.

As diferenças de sequência são calculadas com aritmética sem sinal de 16 bits. O valor delta = seq\_rx - seq\_esperada, reduzido a 16 bits, é considerado antigo ou fora de ordem quando delta \ensuremath{\geq} 32768. O índice na janela também é calculado em 16 bits; posições fora da janela ou já registradas são ignoradas e diagnosticadas. A referência não é utilizada para corrigir o histórico.

O controle usa a sequência recebida também nos descartes porque o injetor tem como alvos somente os bits do payload. Portanto, a implementação assume que o cabeçalho não é corrompido nesse experimento. Essa premissa é necessária: uma sequência pertencente a um quadro com checagem falha não seria confiável em um canal capaz de alterar o cabeçalho.


\par\smallskip\noindent\begin{minipage}{\textwidth}\centering\small
\captionof{table}{Prazos de recepção e seus marcos iniciais}
\begin{tabularx}{\textwidth}{@{}P{.21\linewidth} P{.31\linewidth} Y@{}}\toprule
\textbf{Situação} & \textbf{Prazo} & \textbf{Tratamento}\\\midrule
Corpo ou cópia parcial & 60 ms desde seu marcador & Abandona o bloco e retorna à busca de marcador\\
Referência pendente & 60 ms desde o marcador de dados & Registra auditoria incompleta e invalida a medição\\
Silêncio inicial do enlace & 250 ms desde o primeiro marcador, reiniciado após cada corpo checado & Declara perdida a próxima posição ainda não observada\\
Silêncio persistente & Avanços nominais de 100 ms por posição & Prossegue pela janela sem duplicar classificações\\
\bottomrule\end{tabularx}
\par\smallskip{\raggedright\footnotesize Fonte: Elaboração própria, a partir do RX completo v0.2 (2026).\par}\end{minipage}\par\smallskip

Antes do primeiro marcador de dados, o RX aguarda o início da transmissão e não declara perdas. O timeout de enlace de 250 ms corresponde a aproximadamente 2,48 períodos de 101 ms medidos. O prazo de 60 ms para bloco e auditoria é superior à duração prevista de aproximadamente 45 ms até o fim da referência. Os limites são verificados no loop do firmware.

A passagem de 65535 para zero foi verificada em software; não nas rodadas físicas de 0 a 999.

\FloatBarrier
\clearpage
\section{Materiais e procedimentos}
Foram utilizados três ESP32 clássicos WROOM, um transistor NPN BC547 e dois resistores de 1 k\ensuremath{\Omega}. A saída GPIO17 do TX foi ligada ao GPIO16 do RX por um resistor em série. O coletor do NPN foi conectado ao lado RX da linha e o emissor ao GND comum. A base foi acionada pelo GPIO19 do ESP32 de ruído através do segundo resistor. O injetor observou a linha limpa no GPIO16 e a perturbada no GPIO18.

As placas foram alimentadas individualmente por USB, com os GNDs interligados. O firmware foi compilado para ESP32 Dev Module, CPU de 240 MHz e Arduino-ESP32 3.3.12. A série selecionada empregou o TX preservado, RX completo v0.2, injetor com captura v1.3.1 e coletor CSV v2.


\par\smallskip\noindent\begin{minipage}{\textwidth}\centering\small
\captionof{table}{Modos de ruído empregados nas sete rodadas}
\begin{tabularx}{\textwidth}{@{}c Y@{}}\toprule
\textbf{Modo} & \textbf{Alvos no payload}\\\midrule
0 & Sem pulso; controle de recepção\\
1 & Um bit sorteado entre os 32 bits\\
2 & Dois bits distintos sorteados entre os 32 bits\\
3 & Rajada de 3 a 8 bits contíguos dentro de um byte sorteado\\
\bottomrule\end{tabularx}
\par\smallskip{\raggedright\footnotesize Fonte: Elaboração própria, a partir do injetor e da metodologia experimental (2026).\par}\end{minipage}\par\smallskip

O transistor força somente 1 para 0. Um alvo que já era zero não provoca alteração. A captura de borda em hardware, por MCPWM0, fornece a referência temporal para os pulsos e amostras, sem delayMicroseconds(). Os registros contêm máscara alvo, byte antes, byte depois, alteração efetiva e diagnósticos de sincronismo e atraso. O marcador, LEN, sequência, CHK e cópia não são alvos.

Foram coletados 1000 quadros para SUM8 nos modos 0, 1, 2 e 3, e outros 1000 para CRC8 nos modos 1, 2 e 3. Os payloads e alvos foram sorteados novamente a cada execução. O TX fixa o bit mais significativo de p0 e p1 em 1 para permitir um padrão adicional de demonstração, condição mantida em todas as rodadas.

O coletor salvou simultaneamente os CSVs e logs brutos do RX e do injetor. A conferência incluiu sequência única de 0 a 999, soma das quatro classes, recálculo do resíduo e comparação dos bytes do injetor com a referência e o corpo recebido. Tentativas inválidas não foram incluídas nas tabelas.

\FloatBarrier
\clearpage
\section{Resultados experimentais}
A série contém 7000 registros RX e 11561 registros do injetor. O número de amostras do injetor varia com a quantidade de bytes observados por quadro e não corresponde ao número de quadros. Todas as sete rodadas selecionadas registraram validação de rodada positiva, sem falhas de referência ou divergências nas verificações realizadas.


\par\smallskip\noindent\begin{minipage}{\textwidth}\centering\small
\captionof{table}{Resultados SUM8, com 1000 quadros por modo}
\begin{tabularx}{\textwidth}{@{}Y r r r r r@{}}\toprule
\textbf{Modo} & \textbf{OK} & \textbf{Corrompido} & \textbf{Perdido} & \textbf{Não detectado} & \textbf{Taxa ND}\\\midrule
0 & 1000 & 0 & 0 & 0 & NA\\
1 & 469 & 531 & 0 & 0 & 0,0000\%\\
2 & 213 & 783 & 0 & 4 & 0,5083\%\\
3 & 44 & 956 & 0 & 0 & 0,0000\%\\
\bottomrule\end{tabularx}
\par\smallskip{\raggedright\footnotesize Fonte: Elaboração própria, a partir dos registros experimentais (2026).\par}\end{minipage}\par\smallskip


\par\smallskip\noindent\begin{minipage}{\textwidth}\centering\small
\captionof{table}{Resultados CRC8, com 1000 quadros por modo}
\begin{tabularx}{\textwidth}{@{}Y r r r r r@{}}\toprule
\textbf{Modo} & \textbf{OK} & \textbf{Corrompido} & \textbf{Perdido} & \textbf{Não detectado} & \textbf{Taxa ND}\\\midrule
1 & 443 & 557 & 0 & 0 & 0,0000\%\\
2 & 195 & 805 & 0 & 0 & 0,0000\%\\
3 & 26 & 974 & 0 & 0 & 0,0000\%\\
\bottomrule\end{tabularx}
\par\smallskip{\raggedright\footnotesize Fonte: Elaboração própria, a partir dos registros experimentais (2026).\par}\end{minipage}\par\smallskip

OK identifica um quadro aceito com payload igual à referência. CORROMPIDO identifica um corpo checado e descartado. PERDIDO indica que a checagem não chegou a ser concluída e a falta foi identificada pelo enlace. NÃO DETECTADO identifica uma aceitação com payload diferente da referência válida.

A taxa de não detecção utiliza apenas os quadros com erro de conteúdo:

\begin{equation}T_{\mathrm{ND}}=\frac{N_{\mathrm{ND}}}{N_{\mathrm{C}}+N_{\mathrm{ND}}}.\end{equation}
Para SUM8 no modo 2, a taxa foi 4/(783+4) = 4/787 = 0,005082592, ou 0,5082592\%. Os quatro casos correspondem a 0,4\% de todos os 1000 quadros, fração com denominador diferente. No controle sem erro, a taxa é NA, pois o denominador é zero. Não foram registradas perdas nas sete rodadas regulares.

\FloatBarrier
\clearpage
\subsection{Análise de um erro não detectado}
O quadro 232 da rodada SUM8/modo 2 apresentou duas alterações físicas confirmadas pelos registros do injetor. O corpo original e o recebido, incluindo o detector, foram:

\begin{quote}\small\ttfamily\raggedright Original: 07 E8 00 95 8E B4 99 A1\\
Recebido: 07 E8 00 15 0E B4 99 A1\end{quote}

\par\smallskip\noindent\begin{minipage}{\textwidth}\centering\small
\captionof{table}{Comparação byte a byte do quadro 232}
\begin{tabularx}{\textwidth}{@{}Y r r r r r@{}}\toprule
\textbf{Campo} & \textbf{Original hex} & \textbf{Original decimal} & \textbf{RX hex} & \textbf{RX decimal} & \textbf{Variação}\\\midrule
LEN & 07 & 7 & 07 & 7 & 0\\
seqL & E8 & 232 & E8 & 232 & 0\\
seqH & 00 & 0 & 00 & 0 & 0\\
p0 & 95 & 149 & 15 & 21 & -128\\
p1 & 8E & 142 & 0E & 14 & -128\\
p2 & B4 & 180 & B4 & 180 & 0\\
p3 & 99 & 153 & 99 & 153 & 0\\
CHK & A1 & 161 & A1 & 161 & 0\\
\bottomrule\end{tabularx}
\par\smallskip{\raggedright\footnotesize Fonte: Elaboração própria, a partir dos registros experimentais (2026).\par}\end{minipage}\par\smallskip

A sequência 0x00E8 equivale a 232. A soma dos sete bytes originais é 863, e CHK = -863 mod 256 = 161 = A1. A soma original completa é 1024. No corpo recebido, 95 (149) tornou-se 15 (21), e 8E (142) tornou-se 0E (14). A redução total foi de 256, produzindo soma 768. Como 1024 e 768 são múltiplos de 256, ambos têm resíduo zero.


\par\smallskip\noindent\begin{minipage}{\textwidth}\centering\small
\captionof{table}{Efeito registrado pelo injetor no quadro 232}
\begin{tabularx}{\textwidth}{@{}P{.13\linewidth} P{.16\linewidth} P{.18\linewidth} P{.20\linewidth} Y@{}}\toprule
\textbf{Byte payload} & \textbf{Alvo / alterado} & \textbf{Antes \ensuremath{\rightarrow} depois} & \textbf{Timestamp (\ensuremath{\mu}s)} & \textbf{Mudou / sincronismo}\\\midrule
0 & 80 / 80 & 95 \ensuremath{\rightarrow} 15 & 888739020 & SIM / SIM\\
1 & 80 / 80 & 8E \ensuremath{\rightarrow} 0E & 888742020 & SIM / SIM\\
\bottomrule\end{tabularx}
\par\smallskip{\raggedright\footnotesize Fonte: Elaboração própria, a partir dos registros experimentais (2026).\par}\end{minipage}\par\smallskip

O RX decidiu ACEITO antes de ler a cópia. A auditoria posterior constatou payload diferente e classificou o quadro como NÃO DETECTADO. Os dois registros do injetor concordam com os corpos RX e de referência; os campos de amostras e pulsos tardios foram zero. Os registros integrais estão reproduzidos no Apêndice A.

\FloatBarrier
\clearpage
\subsection{Tempos e overhead}
Os tempos de verificação e os intervalos foram obtidos dos CSVs da série atual. O custo de cálculo no TX foi obtido de um controle anterior sem ruído, com dez quadros por técnica e o mesmo firmware TX. Essas amostras são distintas e não foram agregadas.


\par\smallskip\noindent\begin{minipage}{\textwidth}\centering\small
\captionof{table}{Tempos de cálculo e verificação, em microssegundos}
\begin{tabularx}{\textwidth}{@{}Y r r P{.23\linewidth}@{}}\toprule
\textbf{Medição} & \textbf{SUM8} & \textbf{CRC8} & \textbf{Número de amostras}\\\midrule
TX: média de cálculo & 1,100 & 3,000 & 10 por técnica\\
TX: mínimo / máximo & 1 / 2 & 3 / 3 & Mesmo controle\\
RX: média de verificação & 1,000 & 3,000 & 1000 por rodada atual\\
RX: mínimo / máximo & 1 / 1 & 3 / 3 & 1000 por rodada atual\\
\bottomrule\end{tabularx}
\par\smallskip{\raggedright\footnotesize Fonte: CSV das sete rodadas e capturas do controle TX, identificadas nas referências experimentais.\par}\end{minipage}\par\smallskip


\par\smallskip\noindent\begin{minipage}{\textwidth}\centering\small
\captionof{table}{Intervalos entre inícios de quadro no RX}
\begin{tabularx}{\textwidth}{@{}Y r r@{}}\toprule
\textbf{Coleta} & \textbf{Média (\ensuremath{\mu}s)} & \textbf{Mínimo / máximo (\ensuremath{\mu}s)}\\\midrule
S0 SUM8/0 & 101001,001 & 100999 / 101919\\
S1 SUM8/1 & 101000,000 & 100999 / 101001\\
S2 SUM8/2 & 101000,043 & 100999 / 101022\\
S3 SUM8/3 & 101000,958 & 100999 / 101919\\
C1 CRC8/1 & 101001,001 & 100999 / 101919\\
C2 CRC8/2 & 101001,001 & 100999 / 101919\\
C3 CRC8/3 & 101000,000 & 101000 / 101000\\
\bottomrule\end{tabularx}
\par\smallskip{\raggedright\footnotesize Fonte: Elaboração própria, a partir dos registros experimentais (2026).\par}\end{minipage}\par\smallskip

Cada média de intervalo utiliza 999 diferenças; o zero inicial é excluído. O período observado ficou próximo de 101 ms, justificando o timeout de enlace de 250 ms. As leituras usam esp\_timer\_get\_time() e resolução em microssegundos, sem representar a latência completa do enlace nem medições submicrossegundo.


\par\smallskip\noindent\begin{minipage}{\textwidth}\centering\small
\captionof{table}{Composição dos bytes transmitidos por ciclo}
\begin{tabularx}{\textwidth}{@{}Y r@{}}\toprule
\textbf{Componente} & \textbf{Bytes}\\\midrule
Overhead: marcador 4 + LEN 1 + sequência 2 + CHK 1 & 8\\
Payload & 4\\
Quadro de dados completo & 12\\
Depuração: marcador 4 + cópia do corpo 8 & 12\\
Total de dados e depuração, sem as pausas & 24\\
\bottomrule\end{tabularx}
\par\smallskip{\raggedright\footnotesize Fonte: Elaboração própria, a partir do formato do quadro (2026). A cópia é excluída do overhead do protocolo.\par}\end{minipage}\par\smallskip

\FloatBarrier
\clearpage
\subsection{Perda controlada e recuperação}
Foi executado um ensaio complementar com dez quadros CRC8, RX completo v0.2 e ruído desligado. Após checar e auditar a sequência 1, o RX encerrou sua UART2 por uma janela programada de 350 ms. O TX continuou transmitindo. A ausência de recepção foi induzida por software, separadamente das perturbações elétricas dos experimentos principais.


\par\smallskip\noindent\begin{minipage}{\textwidth}\centering\small
\captionof{table}{Instantes registrados durante a interrupção da recepção}
\begin{tabularx}{\textwidth}{@{}Y r@{}}\toprule
\textbf{Evento} & \textbf{Instante no RX (\ensuremath{\mu}s)}\\\midrule
UART encerrada após auditar seq1 & 707244965\\
Primeira perda declarada & 707481280\\
UART reativada & 707595727\\
Seq5: próxima decisão ACEITO & 707634308\\
\bottomrule\end{tabularx}
\par\smallskip{\raggedright\footnotesize Fonte: Capturas do ensaio de recuperação com RX completo v0.2 (17 out. 2026).\par}\end{minipage}\par\smallskip


\par\smallskip\noindent\begin{minipage}{\textwidth}\centering\small
\captionof{table}{Duração da pausa e tempos de recuperação}
\begin{tabularx}{\textwidth}{@{}Y r r@{}}\toprule
\textbf{Medida} & \textbf{Diferença (\ensuremath{\mu}s)} & \textbf{Resultado}\\\midrule
Janela desligada & 707595727 - 707244965 & 350,762 ms\\
Primeira perda até ACEITO & 707634308 - 707481280 & 153,028 ms\\
Reativação até ACEITO & 707634308 - 707595727 & 38,581 ms\\
\bottomrule\end{tabularx}
\par\smallskip{\raggedright\footnotesize Fonte: Elaboração própria, por diferenças de timestamps da Tabela 11.\par}\end{minipage}\par\smallskip

O resumo apresentou sete OK e três PERDIDO, sem CORROMPIDO ou NÃO DETECTADO. As sequências 2 e 3 foram declaradas perdidas por timeout; a 4 foi identificada pela lacuna ao receber a sequência 5. O quadro 5 foi aceito na checagem, auditado como correto e seguido por recepção até a sequência 9. A captura complementar confirmou rodada válida.

O cronômetro de recuperação preservou a primeira detecção de perda do episódio, sem reiniciar ao detectar a lacuna seguinte. A medida de 153,028 ms termina na próxima aceitação do protocolo; em geral, uma auditoria posterior ainda poderia revelar aceitação indevida. Neste ensaio, a referência confirmou o payload do quadro 5.

As sete verificações desse controle tiveram média de 3,143 \ensuremath{\mu}s, mínimo de 3 \ensuremath{\mu}s e máximo de 4 \ensuremath{\mu}s. Os tempos foram calculados com leituras do mesmo relógio RX, sem subtrair timestamps de placas distintas. O ensaio está documentado por capturas do Monitor Serial e não integra os CSVs das sete rodadas de 1000 quadros.

\FloatBarrier
\clearpage
\section{Discussão e conclusão}
\subsection{Discussão dos resultados}
Os quatro erros não detectados em SUM8/modo 2 exibiram o mesmo mecanismo: apagar dois bits de peso 128 em bytes distintos reduziu a soma em 256. Além da sequência 232, os casos ocorreram nas sequências 239, 530 e 609. Nas duas primeiras, a soma completa passou de 1024 para 768; nas duas últimas, de 768 para 512. A informação perdida pela redução modular explica a aceitação de todos esses quadros alterados.

CRC8 não apresentou erros não detectados nas três rodadas, resultado compatível com sua cobertura para um ou dois bits alterados e para alterações limitadas a um byte. A conclusão se refere às classes ensaiadas: um CRC finito também admite colisões fora delas.

A contagem de descartes isoladamente não mede a superioridade de um detector. As rodadas utilizaram novos sorteios de dados e alvos, e o transistor pode selecionar bits que já eram zero. A comparação utiliza a taxa de aceitação indevida entre quadros efetivamente alterados. Também se deve considerar que p0/p1 possuem o bit mais significativo fixado em 1, de modo que o payload não é uniforme em todos os bits.

Os limites do estudo incluem a restrição do ruído ao payload, a quantização temporal em microssegundos e a dependência do injetor em relação ao core e ao hardware testados. A passagem física da sequência 65535 para zero não foi demonstrada nas rodadas finitas. O ensaio de recuperação mostra retorno após uma interrupção de UART, sem representar todas as condições de perda possíveis.

\subsection{Conclusão}
Foi implementado e avaliado um protocolo de telemetria em que a decisão do RX antecede a leitura da referência. As sete rodadas válidas totalizaram 7000 quadros e permitiram observar quatro aceitações indevidas do SUM8, explicadas e confirmadas byte a byte. O CRC8 apresentou maior tempo de cálculo e verificação nos registros utilizados, acompanhando maior cobertura para as classes de erro ensaiadas.

O RX identificou perdas sem duplicar descartes e retomou a aceitação após a pausa controlada. A escolha do detector deve considerar o modelo de erro e a alteração real do dado, pois nem todo pulso produz corrupção.

\FloatBarrier\clearpage\section*{Referências}\addcontentsline{toc}{section}{Referências}
\begingroup\setlength{\parindent}{0pt}\small SMBUS. \emph{SMBus FAQ: What is the CRC-8 polynomial for SMBus?} [S. l.], s.d. Disponível em: \url{https://smbus.org/faq/faq_main.htm}. Acesso em: 17 out. 2026.\par\endgroup

\begingroup\setlength{\parindent}{0pt}\small KOOPMAN, Philip. \emph{8 Bit CRC Zoo.} Pittsburgh: Carnegie Mellon University, 2024a. Disponível em: \url{https://users.ece.cmu.edu/~koopman/crc/crc8.html}. Acesso em: 17 out. 2026.\par\endgroup

\begingroup\setlength{\parindent}{0pt}\small KOOPMAN, Philip. \emph{CRC Zoo Notes.} Pittsburgh: Carnegie Mellon University, 2024b. Disponível em: \url{https://users.ece.cmu.edu/~koopman/crc/notes.html}. Acesso em: 17 out. 2026.\par\endgroup

\subsection*{Registros experimentais e fontes das medições}
\begingroup\setlength{\parindent}{0pt}\small AURÉLIO, Enzo. \emph{Registros experimentais de telemetria com ESP32.} 17 out. 2026. Conjunto de sete rodadas, com rx.csv, ruido.csv, logs brutos, metadados e revisão por rodada. Arquivos na pasta dados/ do material experimental; identificação no Apêndice A.\par\endgroup

\begingroup\setlength{\parindent}{0pt}\small AURÉLIO, Enzo. \emph{Capturas dos controles de cálculo no TX.} 2026. Dez quadros por técnica, com SUM8 e CRC8. Reproduzidas no documento de evidência Parte 4, páginas 4 e 5. Arquivo: evidencias/Parte\_4\_Tempos\_TX\_e\_Ensaios\_Anteriores.pdf. As medições do antigo RX de ensaio contidas nesse documento não compõem os tempos RX da série atual.\par\endgroup

\begingroup\setlength{\parindent}{0pt}\small AURÉLIO, Enzo. \emph{Capturas do ensaio de recuperação com RX completo v0.2.} 2026. Reproduzidas no documento de evidência Parte 6, página 5, incluindo a confirmação complementar da validade. Arquivo: evidencias/Parte\_6\_Ensaio\_RX\_Completo.pdf.\par\endgroup

Os documentos de evidência preservam as capturas dos instrumentos de registro. As tabelas principais utilizam exclusivamente as sete rodadas identificadas no apêndice. Os controles de cálculo TX e de recuperação mantêm suas amostras e configurações próprias.

\FloatBarrier\clearpage\appendix\section{Registros e rastreabilidade}\setcounter{table}{0}\renewcommand{\thetable}{A.\arabic{table}}\renewcommand{\theHtable}{appendix.\arabic{table}}

\par\smallskip\noindent\begin{minipage}{\textwidth}\centering\small
\captionof{table}{Identificação das sete rodadas selecionadas}
\begin{tabularx}{\textwidth}{@{}l l Y@{}}\toprule
\textbf{ID} & \textbf{Condição} & \textbf{Pasta original}\\\midrule
S0 & SUM8/0 & 20261017\_191422\_664\_SUM8\_modo0\_1000\\
S1 & SUM8/1 & 20261017\_184431\_084\_SUM8\_modo1\_1000\\
S2 & SUM8/2 & 20261017\_185353\_632\_SUM8\_modo2\_1000\\
S3 & SUM8/3 & 20261017\_190634\_547\_SUM8\_modo3\_1000\\
C1 & CRC8/1 & 20261017\_185043\_967\_CRC8\_modo1\_1000\\
C2 & CRC8/2 & 20261017\_185759\_222\_CRC8\_modo2\_1000\\
C3 & CRC8/3 & 20261017\_191057\_739\_CRC8\_modo3\_1000\\
\bottomrule\end{tabularx}
\par\smallskip{\raggedright\footnotesize Fonte: Elaboração própria, a partir dos registros experimentais (2026).\par}\end{minipage}\par\smallskip

Os identificadores S0 a S3 correspondem aos modos SUM8; C1 a C3, aos modos CRC8. As pastas do material experimental foram organizadas como dados/TECNICA\_modoN\_1000. Cada uma contém os seis arquivos originais da coleta, preservados byte a byte. O manifesto SHA256 registra a integridade dessas cópias.

\subsection{Registro RX do quadro 232}
\begin{quote}\small\ttfamily\raggedright seq\_enlace=232; seq\_rx=232; tecnica=SUM8; residuo=0;\\
decisao=ACEITO; status=NAO\_DETECTADO; payload\_confere=NAO;\\
verificacao\_us=1; intervalo\_inicio\_us=101000;\\
corpo\_rx\_hex=07E800150EB499A1;\\
corpo\_tx\_hex=07E800958EB499A1;\\
motivo=CHECADO\_ANTES\_DA\_COPIA.\end{quote}
\subsection{Linhas do injetor para o mesmo quadro}
Campos do CSV, na ordem original: sequência, modo, byte\_payload, máscara alvo, timestamp\_us, antes, depois amostrado, máscara alterada, mudou, sincronismo, atraso máximo, amostras tardias, pulsos tardios e duração do byte.

\begin{quote}\small\ttfamily\raggedright 232,2,0,80,888739020,95,15,80,\\
SIM,SIM,1,00,00,89\end{quote}
\begin{quote}\small\ttfamily\raggedright 232,2,1,80,888742020,8E,0E,80,\\
SIM,SIM,1,00,00,89\end{quote}
As quebras de linha acima servem apenas à diagramação. Nos arquivos originais, cada registro ocupa uma linha. A máscara 0x80 corresponde ao bit 7; em ambos os bytes, o registro informa alteração efetiva e sincronismo válido.


\end{document}
